\chapter{Algorithms}\label{ch:algorithms}

The estimators of Chapter~\ref{ch:estimation} are computed by iterating the
E-step and M-step. This chapter describes the resulting algorithm and its
stopping rule, and it uses the notation of Chapter~\ref{ch:model}.

\section{Expectation maximisation loop}\label{sec:em}

Starting from an initial parameter $\theta^{(0)}$, the algorithm alternates
between computing the responsibilities of Equation~\eqref{eq:resp} and applying
the updates of Equations~\eqref{eq:upd-pi}--\eqref{eq:upd-sigma}. Each
iteration is guaranteed not to decrease the log-likelihood of
Equation~\eqref{eq:loglik}.

\lipsum[27-28]

\section{Initialisation}\label{sec:init}

Good starting values matter. We use $k$-means on the observations to set the
means, and assign each covariance to the empirical covariance of its cluster.
Table~\ref{tab:init} summarises the choices compared in
Chapter~\ref{ch:experiments}.

\begin{table}[ht]
  \centering
  \begin{tabular}{lp{6cm}}
    \toprule
    Strategy & Description \\
    \midrule
    Random & Means drawn uniformly from the data range. \\
    $k$-means & Means from a $k$-means run, covariances from clusters. \\
    Hierarchical & Means from a single-linkage dendrogram cut. \\
    \bottomrule
  \end{tabular}
  \caption{Initialisation strategies.}\label{tab:init}
\end{table}

\lipsum[29-30]

\section{Stopping rule}\label{sec:stopping}

The algorithm stops when the relative change in log-likelihood satisfies
\begin{equation}\label{eq:stop}
  \frac{\ell(\theta^{(t+1)})-\ell(\theta^{(t)})}{|\ell(\theta^{(t)})|} < \tau,
\end{equation}
for a tolerance $\tau$. The default tolerance is $\tau = 10^{-6}$, and the rule
is discussed further in Section~\ref{sec:convergence}.

\lipsum[31-32]

\section{Convergence}\label{sec:convergence}

Because the likelihood is non-decreasing, Equation~\eqref{eq:stop} is a
sensible test. Convergence may nonetheless be slow when components overlap,
and the report returns to this point in Chapter~\ref{ch:discussion}.

\lipsum[33-34]

\section{Further remarks}\label{sec:further-algorithms}

The remarks in this section restate the main point of the chapter in more
general terms, so that the argument can be read without the surrounding
derivations.

\lipsum[99-110]
